%% tcbnestbreak.tex -- documentation for the tcbnestbreak package
%%
%% Copyright (c) 2026 Munehiro Yamamoto <munepixyz@gmail.com>
%% License: LPPL 1.3c or later
%%
%% Build: pdflatex tcbnestbreak.tex (twice)
\documentclass{article}
\usepackage[margin=28mm]{geometry}
\usepackage{xcolor}
\usepackage{tcbnestbreak}
\usepackage[colorlinks, allcolors=blue!50!black]{hyperref}
\newcommand*\tnbversion{v0.2 (2026/11/22)}
\newcommand*\cs[1]{\texttt{\char`\\#1}}
\newcommand*\pkg[1]{\textsf{#1}}
\newcommand*\key[1]{\texttt{#1}}
\newcommand*\meta[1]{$\langle$\textit{#1}$\rangle$}
\tcbset{colframe=black!60, colback=black!3, fonttitle=\bfseries}
\emergencystretch=2em
%% \filler{<n>}: <n> filler sentences
\makeatletter
\newcommand\filler[1]{%
  \count@=0
  \loop\advance\count@ by 1
    This sentence only fills the box, so that the box is long enough to
    be broken.
  \ifnum\count@<#1 \repeat}
\makeatother

\title{The \pkg{tcbnestbreak} package\\
  \large Breakable tcolorboxes inside breakable tcolorboxes}
\author{Munehiro Yamamoto}
\date{\tnbversion}

\begin{document}
\maketitle

\section{Introduction}

\pkg{tcolorbox} typesets the content of a \key{breakable} box as one
box and splits it into the parts for the pages.  A box inside is one
picture in this content and cannot be split, so \pkg{tcolorbox}
ignores \key{breakable} for a box inside another box: a long box
inside a breakable box sticks out of the page, or leaves a large gap
before it.

With the \pkg{tcbnestbreak} package, \key{breakable} works for a box
inside a breakable box, and for a box inside such a box, to any depth.
The box below is a breakable box with a breakable box inside, which
has a lower part and a breakable box of its own.

\begin{tcolorbox}[breakable, title=A breakable box]
The outer box starts here.  \filler{3}

\begin{tcolorbox}[breakable, colback=yellow!15, colframe=blue!50!black,
  title=A breakable box inside]
The inner box starts here.  \filler{6}
\begin{tcolorbox}[breakable, colback=red!8, colframe=red!50!black,
  title=A breakable box inside the inner box]
The innermost box starts here.  \filler{10}
The innermost box ends here.
\end{tcolorbox}
\filler{4}
\tcblower
The lower part of the inner box starts here.  \filler{6}
The inner box ends here.
\end{tcolorbox}

\filler{3}  The outer box ends here.
\end{tcolorbox}

\noindent
Where the outer box is broken, the frame of a box inside is open and
goes on into the padding of the outer box, so that it is seen to
continue on the next page.

\paragraph{Supported formats/engines}
pdf\LaTeX, Xe\LaTeX, Lua\LaTeX{} (with or without Lua\TeX-ja),
p\LaTeX{} and up\LaTeX{} (with dvipdfmx); horizontal (yoko) and
vertical (tate) writing.  The boxes are broken in one run, and no Lua
code is used.  The package requires \pkg{tcolorbox} and loads its
\texttt{breakable} and \texttt{skins} libraries.

\section{Usage}

\begin{verbatim}
\usepackage{tcbnestbreak}

\begin{tcolorbox}[breakable, title=Outer]
  ...
  \begin{tcolorbox}[breakable, title=Inner]
    ...  % broken across pages together with the outer box
  \end{tcolorbox}
  ...
\end{tcolorbox}
\end{verbatim}

\noindent
There is no new key: the package makes \key{breakable} work for a box
that is at the top level of the content of a breakable box, i.e.\
between its paragraphs, also inside a list.  Such a box can have a
title and a lower part (\cs{tcblower}), and it can be defined with
\cs{newtcolorbox}.  A box in a \texttt{minipage}, a \cs{parbox} or a
cell of a \texttt{tabular} is not at the top level: it is an ordinary
unbreakable box, as without the package.

\section{How it works}

\begin{enumerate}
\item \textbf{The inner box flows.}  Its content is typeset indented,
  as the content of a list is (with \cs{parshape}), and put line by
  line into the list of the box around it, with the space of its frame
  above and below.  A \TeX{} mark
  (\cs{marks} of a class of its own) is put at its start, at its end
  and between its upper and its lower part.  \pkg{tcolorbox} then
  splits the outer box between these lines as usual.
\item \textbf{The positions are measured.}  Before a part of the outer
  box is drawn, a copy of the part is taken apart with
  \verb|\vsplit ... to -1sp|, piece by piece.  The heights of the
  pieces and of what is discarded between them (\cs{splitdiscards})
  add up to the position of each mark.
\item \textbf{The frames are drawn.}  For each inner box in the part,
  \pkg{tcolorbox} draws the frame with the options of that box and
  without text: as an unbroken box, or as a first, middle or last part
  of the height that is needed, with the segmentation line if the part
  has it.  The frames are put behind the text of the part.
\end{enumerate}
The text color of an inner box (\key{coltext}, \key{colupper},
\key{collower}) is set again at the top of the next part.

\section{With the \pkg{tcbpagefootnote} package}

The \pkg{tcbpagefootnote} package (v1.4 or later) puts the footnotes
of a box with the key \key{page footnote} at the bottom of the page,
also for each part of a breakable box.  With \pkg{tcbnestbreak}, this
holds for the boxes inside as well: if the outer box and a box inside
have \key{page footnote}, each footnote of the inner box is put on the
page of its mark.  A box inside without the key keeps its footnotes,
below its text.

\section{Limitations}

\begin{itemize}
\item A box flows only at the top level of the content of a breakable
  box or of a box that flows, and not with \key{float},
  \key{sidebyside} or \cs{tcbox}.
\item Where the outer box is broken, the frame of the inner box is
  always open: a part of a box of the \texttt{standard} skin is drawn
  with the skins \texttt{enhancedfirst}, \texttt{enhancedmiddle} and
  \texttt{enhancedlast}, and \key{toprule at break},
  \key{bottomrule at break}, \key{topsep at break} and
  \key{bottomsep at break} of the inner box are not used.
\item The frame is drawn when the part of the outer box is drawn, with
  the options of the inner box applied again.  Colors or styles that
  are defined inside the content of the outer box are not known then.
  \key{height} and the \texttt{enlarge ... by} keys of a box that
  flows are not used.
\item The content is indented like the content of a list.  Material
  that does not follow \cs{parshape} (e.g.\ an \cs{hrule} in vertical
  mode) is not indented.
\item The text color of the box is continued after a break; a
  \cs{color} inside the content is not.
\item The package uses internal macros of \pkg{tcolorbox} (they are
  listed at the top of \texttt{tcbnestbreak.sty}); it is tested with
  \pkg{tcolorbox} 6.9.0 (\TeX{} Live 2025, frozen) and 6.10.0
  (\TeX{} Live 2026).
\end{itemize}

\section{Tests}

\verb|make test| compiles the test files with pdf\LaTeX, Xe\LaTeX,
Lua\LaTeX, p\LaTeX{} and up\LaTeX{} (+ dvipdfmx) and checks the log.
Through the hooks \texttt{tcbnestbreak/flow} and
\texttt{tcbnestbreak/frame}, \texttt{test-common.tex} records in the
\texttt{.aux} file what has flowed and every frame that is drawn; when
the \texttt{.aux} file is read back at \verb|\end{document}|, it
checks that every box has its frames in the right order, on pages that
follow each other and with the segmentation in one of them, and that
the frames together are as high as what has flowed
(\texttt{TNB-RESULT OK}).

\begin{center}
\begin{tabular}{ll}
  \hline
  Test & What it covers\\
  \hline
  \texttt{test-basic.tex}    & boxes inside a box, three levels deep;
                               titles, lower parts, skins,\\
                             & text colors, lists; boxes that do not
                               flow\\
  \texttt{test-plain.tex}    & documents without a box that flows:
                               unchanged by the package\\
  \texttt{test-footnote.tex} & with \pkg{tcbpagefootnote}: the
                               footnotes are on the pages of their
                               marks\\
  \texttt{test-tate.tex}     & \texttt{test-basic.tex} and
                               \texttt{test-footnote.tex} in tate\\
                             & (tarticle, utarticle, ltjtarticle)\\
  \hline
\end{tabular}
\end{center}

\section{History}

\begin{description}
\item[v0.2 (2026/11/22)] A box that flows can have a lower part: the
  segmentation line is drawn, and the text color of the lower part is
  continued.  Vertical writing is tested.  With \pkg{tcbpagefootnote},
  the footnotes of a box inside are put on the pages of their marks.
\item[v0.1 (2026/11/22)] First version.
\end{description}

\section{License}

This package is licensed under the conditions of the \LaTeX{} Project
Public License, either version 1.3c of this license or (at your
option) any later version.  The latest version of this license is in
\url{https://www.latex-project.org/lppl.txt}.

\medskip
\noindent
Copyright \copyright{} 2026 Munehiro Yamamoto

\end{document}
